%% =====================================================================
%%  Unit 14 — What's the weather like?  (Le temps et les saisons)
%%  Pistes : 14.01–14.04 dialogues · 14.05–14.06 vocabulaire · 14.07–14.08 grammaire
%%           14.09 prononciation · 14.10–14.13 exercices (textes en annexe)
%% =====================================================================

\newcommand{\otIllusUnXIV}{\otIllusDuo{sunny/104/86/15/0.10,rain/190/92/20/0.09,winter/122/14/13/0.08,
  March/150/8/11/0.08,cold/196/22/12/0.07,first/100/50/12/0.07}{jack}{What's the weather like?}{lea}{It's raining!}}

\unite[What's the\\weather like?]{14}{What's the weather like?}{Le temps et les saisons}{otIllusUnXIV}

\situation{Épisode 14}{Un lundi matin de novembre. Il pleut à verse sur Manhattan et Léa arrive au
studio trempée de la tête aux pieds. À midi, la conversation passe de la météo… aux anniversaires :
celui de Jack, c'est samedi !}

\begin{objectifs}
\item parler du temps qu'il fait et des saisons ;
\item comprendre une température en degrés Fahrenheit ;
\item dire les mois de l'année et compter avec les nombres ordinaux (\en{first}, \en{second}…) ;
\item donner une date, à l'américaine et à la britannique ;
\item employer \en{in}, \en{on} et \en{at} avec les expressions de temps ;
\item prononcer le son \phon{ɝː} de \en{first}, \en{bird}, \en{work}.
\end{objectifs}

\auprogramme
  {Parler de la météo, donner une date}
  {Les nombres ordinaux\par\vspace{0.6mm}\en{in}, \en{on}, \en{at}}
  {Météo, saisons, dates : 23 expressions, les 12 mois}
  {Le son \phon{ɝː} de \en{first}}
  {Degrés Fahrenheit et jours fériés}

\semaine{casque/otNavy/Dialogues/Écouter et lire les deux scènes,
  bulle/otGreenDark/Vocabulaire/La météo et les mois,
  blocs/otViolet/Grammaire/Ordinaux et prépositions de temps,
  sons/otOrangeDark/Sons et culture/Le son de first et les fêtes,
  crayon/otNavy/Exercices/12 exercices corrigés,
  coche/otGreenDark/Bilan/L'essentiel et le test express}

\clearpage
%% ---------------------------------------------------------------------
\rubrique{ecouter}{La pluie et le beau temps}

\begin{dialogue}{1}{What's the weather like?}{Lundi, 9 heures, au studio. Léa arrive, son parapluie
  retourné à la main.}{14.01}{14.02}
\dit{jack}{Good morning, Léa! Oh, you're \hl{wet}!}{Bonjour, Léa ! Oh, tu es trempée !}
\dit{lea}{Yes! \hl{It's raining} and \hl{it's very windy}. Look at my umbrella!}{Oui ! Il pleut et il y a beaucoup de vent. Regarde mon parapluie !}
\dit{emma}{Welcome to New York \hl{in November}! \hl{What's the weather like} in Lyon?\note{1}}{Bienvenue à New York en novembre ! Quel temps fait-il à Lyon ?}
\dit{lea}{It's gray and \hl{cool}, but it isn't windy like here.}{Il fait gris et frais, mais il n'y a pas de vent comme ici.}
\dit{jack}{My phone says \hl{it's forty degrees} today.}{Mon téléphone dit qu'il fait quarante degrés aujourd'hui.}
\dit{lea}{Forty degrees? But that's very \hl{hot}!}{Quarante degrés ? Mais il fait très chaud !}
\dit{jack}{No, Léa! \hl{Forty degrees Fahrenheit}! That's about four degrees Celsius.\note{2}}{Mais non, Léa ! Quarante degrés Fahrenheit ! Ça fait environ quatre degrés Celsius.}
\dit{emma}{Ha ha! And \hl{in winter}, it's \hl{freezing}. You need a warm coat!\note{3}}{Ha ha ! Et en hiver, il gèle. Il te faut un manteau chaud !}
\end{dialogue}

\begin{comprendre}
\pointnote{1} \en{\textbf{What's the weather like?}} « Quel temps fait-il ? ». Ici, \en{like} ne veut
  pas dire « aimer » : \en{What's it like?} signifie « Comment est-ce ? ». La réponse commence
  toujours par \en{It's…} : \en{It's raining}, \en{It's cold}.
\pointnote{2} \en{\textbf{forty degrees Fahrenheit}} Les Américains mesurent la température en degrés
  Fahrenheit : 40\,°F font environ 4\,°C. Voir la rubrique culture.
\pointnote{3} \en{\textbf{in winter}} On emploie \en{in} devant une saison, un mois ou une année :
  \en{in winter}, \en{in November}, \en{in 2027}.
\end{comprendre}

\begin{dialogue}{2}{When's your birthday?}{Lundi, 13 heures. Pause déjeuner à la cafétéria de
  l'immeuble.}{14.03}{14.04}
\dit{emma}{Jack, \hl{when's your birthday}?\note{4}}{Jack, c'est quand, ton anniversaire ?}
\dit[It's on November twenty-second.]{jack}{It's \hl{on November 22nd}.\note{5}}{C'est le 22 novembre.}
\dit{emma}{But that's this Saturday!}{Mais c'est samedi !}
\dit{jack}{Yes! Come to my place on Saturday, \hl{at seven}.\note{6}}{Oui ! Venez chez moi samedi, à sept heures.}
\dit[Great! My birthday is in March, on March fifteenth.]{lea}{Great! My birthday is \hl{in March}, on March 15th.}{Super ! Mon anniversaire est en mars, le 15 mars.}
\dit{emma}{Mine is \hl{on the first of August}. In summer!}{Le mien, c'est le premier août. En été !}
\dit{jack}{Ha! That's British English. Here, we say \hl{August first}.\note{7}}{Ha ! Ça, c'est de l'anglais britannique. Ici, on dit « August first ».}
\dit{emma}{Whatever, Jack! See you on Saturday!}{Peu importe, Jack ! À samedi !}
\end{dialogue}

\begin{comprendre}
\pointnote{4} \en{\textbf{When's your birthday?}} \en{When's} = \en{When is}, « quand est… ? ».
  \en{birthday} = « anniversaire » (de naissance).
\pointnote{5} \en{\textbf{on November 22nd}} Pour une date, on emploie \en{on} et un nombre
  \textbf{ordinal} : \en{22nd} se lit \en{twenty-second} (« vingt-deuxième »).
\pointnote{6} \en{\textbf{at seven}} On emploie \en{at} devant une heure : \en{at seven}, \en{at
  noon}.
\pointnote{7} \en{\textbf{August first}} Les Américains disent le mois, puis le jour : \en{August
  first}. Les Britanniques disent plutôt \en{the first of August}.
\end{comprendre}

%% ---------------------------------------------------------------------
\rubrique{mots}{Les mots de l'unité}

Voici les 23 expressions de l'unité, puis les douze mois. Pour parler du temps, on commence presque
toujours par \en{It's…} : \en{It's sunny}, \en{It's cold}, \en{It's raining}.

\begin{vocabulaire}{La météo et les saisons}{14.05}
\categorie{Le temps qu'il fait}
\mot{What's the weather like?}{ˌwʌts ðə ˈweðɚ laɪk}{Quel temps fait-il ?}
\mot{It's sunny.}{ɪts ˈsʌni}{Il fait soleil.}
\mot{It's cloudy.}{ɪts ˈklaʊdi}{Le ciel est nuageux.}
\mot{It's raining.}{ɪts ˈreɪnɪŋ}{Il pleut.}
\mot{It's snowing.}{ɪts ˈsnoʊɪŋ}{Il neige.}
\mot{It's windy.}{ɪts ˈwɪndi}{Il y a du vent.}
\mot{It's foggy.}{ɪts ˈfɑːɡi}{Il y a du brouillard.}
\categorie{La température}
\mot{It's hot.}{ɪts ˈhɑːt}{Il fait chaud.}
\mot{It's warm.}{ɪts ˈwɔːrm}{Il fait bon, il fait doux.}
\mot{It's cool.}{ɪts ˈkuːl}{Il fait frais.}
\mot{It's cold.}{ɪts ˈkoʊld}{Il fait froid.}
\mot{It's freezing!}{ɪts ˈfriːzɪŋ}{Il gèle ! Il fait un froid glacial !}
\mot{degrees}{dɪˈɡriːz}{degrés}
\mot{wet}{wet}{mouillé, trempé}
\categorie{Les saisons}
\mot{spring}{sprɪŋ}{le printemps}
\mot{summer}{ˈsʌmɚ}{l'été}
\mot{fall}{fɑːl}{l'automne \registre{américain}}
\mot{autumn}{ˈɔːtəm}{l'automne \registre{britannique}}
\mot{winter}{ˈwɪntɚ}{l'hiver}
\categorie{Les dates}
\mot{When's your birthday?}{ˌwenz jʊr ˈbɝːθdeɪ}{C'est quand, ton anniversaire ?}
\mot[It's on March fifteenth.]{It's on March 15th.}{ɪts ɑːn ˌmɑːrtʃ fɪfˈtiːnθ}{C'est le 15 mars.}
\mot{a holiday}{ə ˈhɑːlədeɪ}{un jour férié, une fête}
\mot{next week}{ˌnekst ˈwiːk}{la semaine prochaine}
\end{vocabulaire}

\begin{grille}{14.06}{Les mois de l'année}{4}
\grcase[January.]{January}{ˈdʒænjueri}{janvier} \grcase[February.]{February}{ˈfebrueri}{février}
\grcase[March.]{March}{mɑːrtʃ}{mars} \grcase[April.]{April}{ˈeɪprəl}{avril}
\grcase[May.]{May}{meɪ}{mai} \grcase[June.]{June}{dʒuːn}{juin}
\grcase[July.]{July}{dʒʊˈlaɪ}{juillet} \grcase[August.]{August}{ˈɑːɡəst}{août}
\grcase[September.]{September}{sepˈtembɚ}{septembre} \grcase[October.]{October}{ɑːkˈtoʊbɚ}{octobre}
\grcase[November.]{November}{noʊˈvembɚ}{novembre} \grcase[December.]{December}{dɪˈsembɚ}{décembre}
\end{grille}

\begin{attention}[Pas de majuscule ? Si !]
En anglais, les mois et les jours prennent toujours une majuscule : \en{in March}, \en{on Monday}.
Les saisons, elles, s'écrivent en minuscules : \en{in spring}, \en{in winter}.
\end{attention}

%% ---------------------------------------------------------------------
\rubrique{grammaire}{Les dates et les prépositions de temps}

\otIntertitre[otViolet]{1 · Les nombres ordinaux}
Pour classer (« premier », « deuxième »…) et pour les dates, on ajoute \en{-th} au nombre :
\en{four} → \en{fourth}. Quelques formes sont irrégulières : apprenez-les par cœur.

\begin{grille}{14.07}{Les nombres ordinaux}{5}
\grcase[First.]{1st}{fɝːst}{first} \grcase[Second.]{2nd}{ˈsekənd}{second} \grcase[Third.]{3rd}{θɝːd}{third}
\grcase[Fourth.]{4th}{fɔːrθ}{fourth} \grcase[Fifth.]{5th}{fɪfθ}{fifth}
\grcase[The sixth day.]{6th}{sɪksθ}{sixth} \grcase[Seventh.]{7th}{ˈsevənθ}{seventh} \grcase[Eighth.]{8th}{eɪtθ}{eighth}
\grcase[Ninth.]{9th}{naɪnθ}{ninth} \grcase[Tenth.]{10th}{tenθ}{tenth}
\grcase[Twelfth.]{12th}{twelfθ}{twelfth} \grcase[Twentieth.]{20th}{ˈtwentiəθ}{twentieth}
\grcase[Twenty-first.]{21st}{ˌtwentiˈfɝːst}{twenty-first} \grcase[Thirtieth.]{30th}{ˈθɝːtiəθ}{thirtieth}
\grcase[Thirty-first.]{31st}{ˌθɝːtiˈfɝːst}{thirty-first}
\end{grille}

\begin{regle}[Les formes à retenir]
\begin{etapes}
\item Trois formes à part : \en{first} (1\textsuperscript{st}), \en{second} (2\textsuperscript{nd}),
  \en{third} (3\textsuperscript{rd}), et donc \en{twenty-first}, \en{thirty-second}…
\item Orthographe modifiée : \en{fifth}, \en{eighth}, \en{ninth}, \en{twelfth} ; \en{twenty} →
  \en{twentieth}.
\item En abrégé, on écrit les deux dernières lettres : \en{1st}, \en{2nd}, \en{3rd}, \en{4th}, \en{22nd}.
\end{etapes}
\end{regle}

\otIntertitre[otViolet]{2 · Dire la date}
\begin{center}\otTableau
\begin{tabular}{@{}l@{\hspace{8mm}}>{\enOn}l<{\enOff}@{\hspace{8mm}}>{\enOn}l<{\enOff}@{}}
 & \otTH{\en{American English}} & \otTH{\en{British English}}\\\midrule
on écrit & March 15th · 3/15/2027 & 15th March · 15/3/2027\\
on dit   & March fifteenth & the fifteenth of March\\
l'année  & twenty twenty-seven & twenty twenty-seven\\\bottomrule
\end{tabular}\end{center}

\otIntertitre[otViolet]{3 · \en{in}, \en{on} ou \en{at} ?}
\begin{center}\otTableau
\begin{tabular}{@{}>{\enOn\otSemi}l<{\enOff}@{\hspace{8mm}}p{58mm}@{\hspace{8mm}}>{\enOn\raggedright\arraybackslash}p{74mm}<{\enOff}@{}}
\otTH{Préposition} & \otTH{Devant\ldots} & \otTH{Exemples}\\\midrule
at & une heure, un moment précis & at seven, at noon, at night\\
on & un jour, une date & on Monday, on March 15th, on my birthday\\
in & un mois, une saison, une année, une partie de la journée & in May, in winter, in 2027, in the morning\\\bottomrule
\end{tabular}\end{center}

\blocaudio{14.08}{Écoutez et répétez}{Six phrases avec une date ou une préposition de temps.}

\begin{attention}[Les pièges des francophones]
\fauxjuste{It's raining in Monday.}{It's raining on Monday.}
\fauxjuste{My birthday is the 15 March.}{My birthday is on March 15th.}
\fauxjuste{It makes cold.}{It's cold.}
\fauxjuste{in the night}{at night}
\fauxjuste{in january}{in January}
\end{attention}

%% ---------------------------------------------------------------------
\rubrique[Le son de first]{sons}{Le son \phon{ɝː} de \en{first}}

Le son \phon{ɝː} n'existe pas en français. Partez du « eu » de \textit{peur}, puis reculez et
recourbez la langue vers le haut, sans toucher le palais : c'est le \en{r} américain, qui « colore » la
voyelle. Un seul son, quelle que soit l'orthographe : \en{fi\textbf{r}st}, \en{bi\textbf{r}d},
\en{Thu\textbf{r}sday}, \en{he\textbf{r}}, \en{wo\textbf{r}k}, \en{ea\textbf{r}ly}. Ne le confondez pas
avec le \phon{ɔːr} de \en{four}, \en{short}, \en{born}, où les lèvres s'arrondissent.

\begin{paires}{14.09}{Écoutez chaque paire et répétez : \phon{ɝː} (lèvres détendues) puis \phon{ɔːr}
(lèvres arrondies).}
\sonpaire{sh\otH{ir}t}{ʃɝːt}{une chemise}{sh\otH{or}t}{ʃɔːrt}{court}
\sonpaire{b\otH{ir}d}{bɝːd}{un oiseau}{b\otH{oar}d}{bɔːrd}{une planche, un tableau}
\sonpaire{b\otH{ur}n}{bɝːn}{brûler}{b\otH{or}n}{bɔːrn}{né}
\sonpaire{t\otH{ur}n}{tɝːn}{tourner}{t\otH{or}n}{tɔːrn}{déchiré}
\sonpaire{w\otH{or}d}{wɝːd}{un mot}{w\otH{ar}d}{wɔːrd}{un service (d'hôpital)}
\sonpaire{f\otH{ir}m}{fɝːm}{ferme, solide}{f\otH{or}m}{fɔːrm}{un formulaire}
\end{paires}

\begin{astuce}[\en{work} ou \en{walk} ?]
Le piège classique : \en{work} \phon{wɝːk} (travailler) contient le son \phon{ɝː}, alors que \en{walk}
\phon{wɑːk} (marcher) se prononce sans \en{r} ni \en{l} ! Le \en{l} de \en{walk}, \en{talk}, \en{half} est
muet.
\end{astuce}

%% ---------------------------------------------------------------------
\rubrique{culture}{Degrés Fahrenheit et jours fériés}

\begin{culture}[Fahrenheit ou Celsius ?]
Les États-Unis mesurent la température en degrés Fahrenheit (°F). Quelques repères : l'eau gèle à
32\,°F (0\,°C) ; 50\,°F = 10\,°C (frais) ; 68\,°F = 20\,°C (agréable) ; 86\,°F = 30\,°C (chaud) ; une
fièvre commence vers 100\,°F (38\,°C). Pour estimer de tête : retirez 30, puis divisez par 2. Ainsi,
70\,°F ≈ (70 − 30) ÷ 2 = 20\,°C.
\end{culture}

\begin{paire}
\begin{culture}[Les dates à l'américaine]
Aux États-Unis, le mois passe \textbf{avant} le jour : \en{7/4} est le 4 juillet, et non le 7 avril !
Sur un formulaire, \en{03/15/2027} signifie donc 15 mars 2027. Prudence avec les dates de réservation.
\end{culture}
\begin{culture}[Les grands jours fériés]
\en{New Year's Day} (1\textsuperscript{er} janvier), \en{Memorial Day} (fin mai), \en{Independence Day}
(4 juillet : feux d'artifice et barbecues), \en{Labor Day} (début septembre), \en{Thanksgiving} (fin
novembre) et \en{Christmas} (25 décembre).
\end{culture}
\end{paire}

%% ---------------------------------------------------------------------
\rubrique{pratique}{Exercices}

\exercice[lire,ecrire,niveau=1]{Quel temps fait-il ?}{Regardez chaque image et écrivez la phrase qui
convient, avec \en{It's…}.}
\begin{tcbraster}[raster columns=3,raster equal height=rows,raster column skip=3mm,raster row skip=3mm,raster before skip=1mm]
\otCarteObjet{soleil}{otOrange}{1}{It's sunny.}
\otCarteObjet{pluie}{otNavy}{1}{It's raining.}
\otCarteObjet{neige}{otTeal}{1}{It's snowing.}
\otCarteObjet{nuage}{otMuted}{1}{It's cloudy.}
\otCarteObjet{vent}{otViolet}{1}{It's windy.}
\otCarteObjet{brouillard}{otMuted}{1}{It's foggy.}
\end{tcbraster}

\exercice[ecrire,niveau=1]{Les nombres ordinaux}{Écrivez chaque nombre ordinal en lettres.}
\begin{anglais}
\begin{questions}[4]
\item 1st → \trou[20mm]{first}
\item 2nd → \trou[20mm]{second}
\item 3rd → \trou[20mm]{third}
\item 5th → \trou[20mm]{fifth}
\item 9th → \trou[20mm]{ninth}
\item 12th → \trou[20mm]{twelfth}
\item 20th → \trou[20mm]{twentieth}
\item 31st → \trou[20mm]{thirty-first}
\end{questions}
\end{anglais}

\exercice[ecrire,niveau=1]{\en{in}, \en{on} ou \en{at} ?}{Complétez avec la bonne préposition.}
\begin{anglais}
\begin{questions}[2]
\item My birthday is \trou[12mm]{on} March 15th.
\item It's cold \trou[12mm]{in} winter.
\item The concert is \trou[12mm]{at} eight o'clock.
\item I don't work \trou[12mm]{on} Sundays.
\item We have a holiday \trou[12mm]{in} July.
\item I never go out \trou[12mm]{at} night.
\item I read the news \trou[12mm]{in} the morning.
\item See you \trou[12mm]{on} Saturday!
\end{questions}
\end{anglais}

\exercice[lire,niveau=2]{La météo de la semaine}{Voici les prévisions pour New York. Indiquez si
chaque phrase est vraie (V) ou fausse (F).}
\otPrevision{Monday/pluie/otNavy/45/7,Tuesday/nuage/otMuted/48/9,Wednesday/soleil/otOrange/52/11,Thursday/vent/otViolet/41/5,Friday/neige/otTeal/30/−1}
\begin{anglais}
\begin{questions}[2]
\item It's sunny on Wednesday. \hfill\trou[9mm]{V}
\item It's snowing on Monday. \hfill\trou[9mm]{F}
\item On Friday, it's freezing. \hfill\trou[9mm]{V}
\item It's windy on Thursday. \hfill\trou[9mm]{V}
\item On Tuesday, it's very hot. \hfill\trou[9mm]{F}
\item It's warm on Friday. \hfill\trou[9mm]{F}
\end{questions}
\end{anglais}

\exercice[ecrire,niveau=2]{Les anniversaires}{Ces dates sont écrites à l'américaine (mois/jour).
Écrivez-les en toutes lettres, comme on les dit aux États-Unis.}
\begin{questions}
\item Jack : \enb{11/22} \hfill \enb{\trou[60mm]{November twenty-second}}
\item Emma : \enb{8/1} \hfill \enb{\trou[60mm]{August first}}
\item Sam : \enb{6/3} \hfill \enb{\trou[60mm]{June third}}
\item Mrs.\ Lopez : \enb{9/12} \hfill \enb{\trou[60mm]{September twelfth}}
\item Léa : \enb{3/15} \hfill \enb{\trou[60mm]{March fifteenth}}
\item Maya : \enb{4/30} \hfill \enb{\trou[60mm]{April thirtieth}}
\end{questions}

\exercice[lire,niveau=1]{Les fêtes de l'année}{Associez chaque fête à sa date : écrivez la lettre.}
\begin{anglais}
\noindent\begin{minipage}[t]{0.48\linewidth}
\begin{questions}[1][english]
\item New Year's Day \hfill\trou[9mm]{c}
\item Valentine's Day \hfill\trou[9mm]{e}
\item Independence Day \hfill\trou[9mm]{a}
\item Halloween \hfill\trou[9mm]{f}
\item Thanksgiving \hfill\trou[9mm]{b}
\item Christmas \hfill\trou[9mm]{d}
\end{questions}
\end{minipage}\hfill
\begin{minipage}[t]{0.46\linewidth}
\begin{questions}[1][english]
\item[a.] July 4th
\item[b.] the fourth Thursday in November
\item[c.] January 1st
\item[d.] December 25th
\item[e.] February 14th
\item[f.] October 31st
\end{questions}
\end{minipage}
\end{anglais}

\exercice[ecouter,niveau=1,piste=14.10]{\phon{ɝː} ou \phon{ɔːr} ?}{Écoutez. Pour chaque numéro, la
voix dit un seul des deux mots, deux fois. Cochez le mot entendu.}
\begin{questions}[3]
\item \choix{shirt}\quad\choix{short}
\item \choix{bird}\quad\choix{board}
\item \choix{burn}\quad\choix{born}
\item \choix{turn}\quad\choix{torn}
\item \choix{word}\quad\choix{ward}
\item \choix{firm}\quad\choix{form}
\end{questions}

\exercice[ecouter,ecrire,niveau=2,piste=14.11]{La météo à la radio}{Écoutez le bulletin météo et
complétez le tableau.}
\begin{center}\otTableau
\begin{tabular}{@{}>{\enOn}l<{\enOff}@{\hspace{14mm}}>{\enOn}l<{\enOff}@{\hspace{14mm}}>{\enOn}l<{\enOff}@{}}
\otTH{City} & \otTH{Weather} & \otTH{Temperature}\\\midrule
New York & It's \trou[24mm]{cloudy}. & \trou[14mm]{50} °F\\[0.8mm]
Boston   & It's \trou[24mm]{raining}. & \trou[14mm]{45} °F\\[0.8mm]
Chicago  & It's \trou[24mm]{windy}. & \trou[14mm]{38} °F\\[0.8mm]
Miami    & It's \trou[24mm]{sunny}. & \trou[14mm]{82} °F\\\bottomrule
\end{tabular}\end{center}

\exercice[ecouter,ecrire,niveau=2,piste=14.12]{Dictée à trous}{Sam et Léa parlent de leurs
anniversaires. Écoutez et complétez.}
\begin{anglais}
\noindent\begin{tabular}{@{}l@{\hspace{6mm}}l@{}}
\otNomDit{sam} & Léa, \trou{when's} your birthday?\\[1.8mm]
\otNomDit{lea} & It's on \trou{March} 15th. And \trou{yours}?\\[1.8mm]
\otNomDit{sam} & It's in \trou{summer}, on \trou{June} 3rd.\\[1.8mm]
\otNomDit{lea} & What's the \trou{weather} like in Atlanta in June?\\[1.8mm]
\otNomDit{sam} & It's very \trou{hot} and sunny! It's ninety degrees!\\[1.8mm]
\otNomDit{lea} & In Lyon, it's \trou{warm} in June, but not hot.\\
\end{tabular}
\end{anglais}

\exercice[ecrire,niveau=2]{Traduisez}{Traduisez en anglais.}
\begin{questions}[1][french]
\item Quel temps fait-il ? — Il pleut et il fait froid.\lignes{1}
\item C'est quand, ton anniversaire ? — C'est le 3 juin.\lignes{1}
\item Il neige beaucoup en hiver.\lignes{1}
\item Il fait 30 degrés Celsius : il fait très chaud !\lignes{1}
\item Le concert est samedi, à huit heures.\lignes{1}
\item Je n'aime pas le vent.\lignes{1}
\item Noël, c'est le 25 décembre.\lignes{1}
\item Je travaille le matin.\lignes{1}
\end{questions}

\exercice[ecrire,niveau=3]{Ma saison préférée}{En cinq phrases, décrivez votre saison préférée : quels
mois, quel temps il fait, ce que vous aimez faire. Donnez aussi la date de votre anniversaire.}
\lignes{5}

\exercice[ecouter,parler,niveau=3,piste=14.13]{À vous de parler !}{Emma vous pose des questions.
Répondez \textbf{à voix haute} pendant chaque pause.}
\begin{center}\otTableau
\begin{tabular}{@{}>{\otSemi\fontsize{8.6}{10}\selectfont\color{otNavy}}l@{\hspace{3mm}}>{\enOn}p{70mm}<{\enOff}@{\hspace{6mm}}p{78mm}@{}}
 & \otTH{Emma dit\ldots} & \otTH{Vous\ldots}\\\midrule
1. & What's the weather like today? & décrivez le temps qu'il fait.\\
2. & Is it cold in your town in winter? & répondez et précisez.\\
3. & What's your favorite season? & donnez votre saison préférée.\\
4. & When's your birthday? & donnez la date.\\
5. & And what's the date today? & donnez la date du jour.\\\bottomrule
\end{tabular}\end{center}

%% ---------------------------------------------------------------------
\rubrique{bilan}{Faites le point}

\begin{essentiel}
\begin{minipage}[t]{0.54\linewidth}
\otEssTitre{Les phrases clés}
\renewcommand{\arraystretch}{1.32}
\begin{tabular}{@{}>{\enOn\otSemi\fontsize{9.4}{11}\selectfont}l<{\enOff}@{\hspace{3mm}}>{\fontsize{9.2}{11}\selectfont\color{white!78!otNavy}}l@{}}
What's the weather like?   & Quel temps fait-il ?\\
It's raining and it's cold. & Il pleut et il fait froid.\\
It's forty degrees.        & Il fait quarante degrés (°F).\\
When's your birthday?      & C'est quand, ton anniversaire ?\\
It's on March 15th.        & C'est le 15 mars.\\
See you on Saturday!       & À samedi !\\
\end{tabular}
\end{minipage}\hfill
\begin{minipage}[t]{0.41\linewidth}
\otEssTitre{Les ordinaux}
{\fontsize{9.2}{11.6}\selectfont \en{first}, \en{second}, \en{third}, puis \en{-th} : \en{fourth},
\en{fifth}, \en{twelfth}, \en{twentieth}\par}
\otEssTitre{\en{at} · \en{on} · \en{in}}
{\fontsize{9.2}{11.6}\selectfont \en{at seven} · \en{on Monday}, \en{on March 15th} · \en{in May},
\en{in winter}, \en{in the morning}\par}
\otEssTitre{Prononciation}
{\fontsize{9.2}{11.6}\selectfont \en{first}, \en{bird}, \en{work} \phon{ɝː} ≠ \en{four}, \en{short}
\phon{ɔːr}\par}
\end{minipage}
\end{essentiel}

\begin{jesais}
\item décrire le temps qu'il fait et nommer les saisons ;
\item comprendre une température en degrés Fahrenheit ;
\item dire les mois et les nombres ordinaux ;
\item donner une date, à l'américaine et à la britannique ;
\item choisir entre \en{at}, \en{on} et \en{in} ;
\item prononcer \phon{ɝː} et le distinguer de \phon{ɔːr}.
\end{jesais}

\clearpage
\begin{testexpress}
\qcm[c]{« Quel temps fait-il ? » se dit :}{\en{How is the weather like?}}{\en{What weather does it make?}}{\en{What's the weather like?}}
\qcm[a]{« Il pleut » se dit :}{\en{It's raining.}}{\en{It rains now.}}{\en{It makes rain.}}
\qcm[b]{« Il fait froid » se dit :}{\en{It makes cold.}}{\en{It's cold.}}{\en{It has cold.}}
\qcm[b]{L'ordinal de \en{five} est :}{\en{fiveth}}{\en{fifth}}{\en{fivth}}
\qcm[a]{« le 22 novembre » s'écrit, à l'américaine :}{\en{November 22nd}}{\en{22th November}}{\en{November 22st}}
\qcm[c]{« en hiver » se dit :}{\en{on winter}}{\en{at winter}}{\en{in winter}}
\qcm[a]{« samedi » (rendez-vous) se dit :}{\en{on Saturday}}{\en{in Saturday}}{\en{at Saturday}}
\qcm[c]{Aux États-Unis, \en{7/4} désigne :}{le 7 avril}{le 7 juillet}{le 4 juillet}
\qcm[b]{68\,°F correspondent environ à :}{68\,°C}{20\,°C}{0\,°C}
\qcm[a]{Dans \en{first}, on entend :}{\phon{ɝː}, comme dans \en{bird}}{\phon{ɔːr}, comme dans \en{four}}{\phon{iː}, comme dans \en{see}}
\end{testexpress}

\noindent{\fontsize{9.6}{12}\selectfont\color{otMuted}Niveau visé (CECRL, niveau A1) : « peut parler
du temps qu'il fait et donner une date en termes simples ». Les corrigés de l'unité sont
page~\pageref{cor-u14-1} ; les textes des pistes d'exercices, page~\pageref{trans-14.10}.\par}
